\documentclass[11pt,a4paper]{article}

% ---------------------------------------------------------
% PACKAGES
% ---------------------------------------------------------
\usepackage[utf8]{inputenc}
\usepackage[T1]{fontenc}
\usepackage{lmodern}
\usepackage{amsmath, amssymb, amsfonts}
\usepackage{bm}
\usepackage{graphicx}
\usepackage{hyperref}
\usepackage{authblk}
\usepackage{geometry}
\usepackage{cite}
\usepackage{physics}
\usepackage{tensor}
\usepackage{enumitem}

\geometry{margin=1in}

% ---------------------------------------------------------
% TITLE
% ---------------------------------------------------------
\title{\textbf{Companion XXI: Reionization and the Thermal History in the Relaxation‑Driven Cyclic Cosmology Ionizing Sources, IGM Heating, and the RDCC Photon–Baryon Connection}}

\author[1]{Michael Lehmann}
\affil[1]{\small Independent Researcher, Relaxation-Driven Cyclic Cosmology (RDCC) Framework\\
\texttt{mi.lehmann@gmx.de}}

\date{\small \today}

% ---------------------------------------------------------
% DOCUMENT START
% ---------------------------------------------------------
\begin{document}

\maketitle

\begin{abstract}
The epoch of reionization (EoR) marks the transition from a neutral to an ionized Universe
and encodes essential information about early structure formation, the first stars and galaxies,
and the thermal history of the intergalactic medium (IGM). In the standard $\Lambda$CDM
model, reionization is driven by early star formation and regulated by feedback, radiative
transfer, and the growth of small-scale structure. In the Relaxation-Driven Cyclic Cosmology
(RDCC), however, the relaxation-induced modification of gravitational dynamics directly
affects the formation of early halos, the production of ionizing photons, and the thermal
evolution of the IGM.

This Companion develops the first RDCC reionization framework. We derive how
relaxation-modified collapse thresholds, scale-dependent growth, and baryonic suppression
alter the abundance of ionizing sources, the escape fraction of UV photons, and the heating
and ionization history of the IGM. We show that RDCC predicts a delayed onset of
reionization, a smoother ionization morphology, reduced small-scale clumping, and a distinct
21-cm signal with a characteristic transition scale tied to the relaxation scale $k_{\rm rel}$.

Building on the baryonic physics developed in Companion~XX and the non-linear structure
formation analysis of Companion~XIX, we construct an RDCC-adjusted reionization model
and derive its impact on the ionization fraction, optical depth, IGM temperature, and 21-cm
brightness temperature. The resulting predictions are consistent with current constraints from
Planck, HERA, LOFAR, and JWST, and provide clear targets for upcoming surveys such as
SKA and the Roman Space Telescope.

This Companion establishes cosmic reionization as a powerful observational probe of the
relaxation dynamics of RDCC and extends the framework into the earliest luminous epoch of
the Universe.
\end{abstract}

% ---------------------------------------------------------
% SECTION 1 — RDCC-MODIFIED IONIZING EMISSIVITY AND ESCAPE FRACTION
% ---------------------------------------------------------
\section{RDCC-Modified Ionizing Emissivity and Escape Fraction}

The ionizing emissivity and the escape fraction of UV photons are the two key ingredients
that determine the timing and morphology of cosmic reionization. In the standard
$\Lambda$CDM framework, both quantities are governed by the abundance of low-mass halos,
the star formation efficiency, and the structure of the interstellar medium (ISM). In RDCC,
the relaxation-modified gravitational dynamics alter each of these components.

Relaxation affects reionization through:
\begin{itemize}
\item suppressed early halo formation (Companion~XIX),
\item reduced star formation efficiency (Companion~XX),
\item enhanced escape fractions in low-mass halos,
\item reduced clumping and smoother gas distributions,
\item a characteristic scale dependence tied to $k_{\mathrm{rel}}$.
\end{itemize}

This section derives the RDCC-modified ionizing emissivity, photon production efficiency,
and escape fraction.

% ---------------------------------------------------------
\subsection{Ionizing Emissivity in RDCC}

The comoving ionizing emissivity is defined as
\begin{equation}
\dot{n}_{\gamma}(z)
=
\int dM\,
\dot{N}_{\gamma}(M,z)\,
n(M,z),
\end{equation}
where $\dot{N}_{\gamma}$ is the ionizing photon production rate of a halo of mass $M$.

In RDCC, the emissivity is modified by:
\begin{itemize}
\item the suppressed halo mass function (Companion~XIX),
\item reduced star formation efficiency (Companion~XX),
\item delayed collapse of low-mass halos,
\item modified gas inflow and cooling rates.
\end{itemize}

We write the RDCC emissivity as
\begin{equation}
\dot{n}_{\gamma,\mathrm{RDCC}}
=
\dot{n}_{\gamma,\Lambda\mathrm{CDM}}
\left[
1
-
\chi_{\gamma}
\left(
\frac{k}{k_{\mathrm{rel}}}
\right)^{\alpha_{\gamma}}
\right],
\end{equation}
with $\chi_{\gamma} \sim 0.1$.

Consequences:
\begin{itemize}
\item suppressed emissivity at $z \gtrsim 10$,
\item delayed onset of reionization,
\item smoother rise of the ionizing background.
\end{itemize}

% ---------------------------------------------------------
\subsection{Photon Production Efficiency}

The ionizing photon production rate is
\begin{equation}
\dot{N}_{\gamma}
=
\zeta_{\mathrm{ion}}
\dot{M}_{\star},
\end{equation}
where $\zeta_{\mathrm{ion}}$ is the ionizing photon production efficiency.

In RDCC, the star formation rate is reduced (Companion~XX):
\begin{equation}
\dot{M}_{\star,\mathrm{RDCC}}
=
\dot{M}_{\star,\Lambda\mathrm{CDM}}
\left[
1
-
\chi_{\star}
\left(
\frac{M_{\mathrm{rel}}}{M}
\right)^{\beta_{\star}}
\right],
\end{equation}
with $\chi_{\star} \sim 0.1$.

Thus,
\begin{equation}
\zeta_{\mathrm{ion,RDCC}}
=
\zeta_{\mathrm{ion},\Lambda\mathrm{CDM}}
\left[
1
-
\chi_{\mathrm{ion}}
\left(
\frac{M_{\mathrm{rel}}}{M}
\right)^{\beta_{\mathrm{ion}}}
\right].
\end{equation}

Consequences:
\begin{itemize}
\item fewer ionizing photons per halo at early times,
\item reduced contribution from dwarf galaxies,
\item delayed buildup of the UV background.
\end{itemize}

% ---------------------------------------------------------
\subsection{Escape Fraction in RDCC}

The escape fraction $f_{\mathrm{esc}}$ is the fraction of ionizing photons that escape the ISM and
reach the IGM.

In RDCC, several effects modify $f_{\mathrm{esc}}$:
\begin{itemize}
\item shallower gravitational potentials reduce gas densities,
\item reduced star formation lowers dust production,
\item enhanced SN-driven outflows create low-density channels,
\item relaxation-induced vorticity disrupts cold gas inflow.
\end{itemize}

We model the RDCC escape fraction as
\begin{equation}
f_{\mathrm{esc,RDCC}}
=
f_{\mathrm{esc},\Lambda\mathrm{CDM}}
\left[
1
+
\chi_{\mathrm{esc}}
\left(
\frac{M_{\mathrm{rel}}}{M}
\right)^{\beta_{\mathrm{esc}}}
\right],
\end{equation}
with $\chi_{\mathrm{esc}} \sim 0.1$.

Consequences:
\begin{itemize}
\item higher escape fractions in low-mass halos,
\item reduced sensitivity to ISM clumping,
\item smoother ionization fronts.
\end{itemize}

% ---------------------------------------------------------
\subsection{Effective Ionizing Efficiency}

The effective ionizing efficiency is
\begin{equation}
\zeta
=
\zeta_{\mathrm{ion}}
f_{\mathrm{esc}}
\epsilon_{\star}.
\end{equation}

In RDCC:
\begin{equation}
\zeta_{\mathrm{RDCC}}
=
\zeta_{\Lambda\mathrm{CDM}}
\left[
1
-
\chi_{\star}
+
\chi_{\mathrm{esc}}
\right]
\left(
\frac{M_{\mathrm{rel}}}{M}
\right)^{\beta_{\zeta}}.
\end{equation}

Thus:
\begin{itemize}
\item reduced SFR dominates at high redshift,
\item enhanced escape fraction partially compensates,
\item the net effect is a delayed but more efficient late-time reionization.
\end{itemize}

% ---------------------------------------------------------
\subsection{Summary}

RDCC modifies ionizing emissivity and escape fractions through:
\begin{itemize}
\item suppressed early star formation,
\item reduced photon production efficiency,
\item enhanced escape fractions in low-mass halos,
\item and a characteristic scale dependence tied to $k_{\mathrm{rel}}$.
\end{itemize}

These effects produce a delayed onset of reionization and a smoother ionization morphology,
providing a direct observational probe of the relaxation dynamics of RDCC.

% ---------------------------------------------------------
% SECTION 2 — RDCC REIONIZATION HISTORY AND OPTICAL DEPTH
% ---------------------------------------------------------
\section{RDCC Reionization History and Optical Depth}

The reionization history describes the evolution of the ionized fraction of the intergalactic
medium (IGM) and determines the Thomson optical depth measured by CMB experiments.
In the standard $\Lambda$CDM model, reionization is driven by early star formation in
low-mass halos and proceeds rapidly once the ionizing emissivity exceeds the recombination
rate. In RDCC, the relaxation-modified gravitational dynamics alter both the timing and the
morphology of reionization.

Relaxation affects reionization through:
\begin{itemize}
\item reduced early ionizing emissivity (Section~1),
\item delayed collapse of low-mass halos (Companion~XIX),
\item reduced clumping factors due to smoother small-scale structure,
\item enhanced escape fractions in low-mass halos,
\item a characteristic scale dependence tied to $k_{\mathrm{rel}}$.
\end{itemize}

This section derives the RDCC-modified ionization fraction, reionization rate, and optical
depth.

% ---------------------------------------------------------
\subsection{Ionization Fraction Evolution}

The volume-averaged ionized fraction evolves according to
\begin{equation}
\frac{dQ_{\mathrm{HII}}}{dt}
=
\frac{\dot{n}_{\gamma}}{\bar{n}_{\mathrm{H}}}
-
\frac{Q_{\mathrm{HII}}}{t_{\mathrm{rec}}},
\end{equation}
where $\dot{n}_{\gamma}$ is the ionizing emissivity and $t_{\mathrm{rec}}$ is the recombination time.

In RDCC, both terms are modified:
\begin{itemize}
\item $\dot{n}_{\gamma}$ is reduced due to suppressed early star formation (Section~1),
\item $t_{\mathrm{rec}}$ increases because reduced small-scale structure lowers the clumping factor.
\end{itemize}

We write the RDCC recombination time as
\begin{equation}
t_{\mathrm{rec,RDCC}}
=
t_{\mathrm{rec},\Lambda\mathrm{CDM}}
\left[
1
+
\chi_{\mathrm{rec}}
\left(
\frac{k}{k_{\mathrm{rel}}}
\right)^{\alpha_{\mathrm{rec}}}
\right],
\end{equation}
with $\chi_{\mathrm{rec}} \sim 0.1$.

Thus, the RDCC ionization fraction evolves as
\begin{equation}
\frac{dQ_{\mathrm{HII,RDCC}}}{dt}
=
\frac{\dot{n}_{\gamma,\mathrm{RDCC}}}{\bar{n}_{\mathrm{H}}}
-
\frac{Q_{\mathrm{HII,RDCC}}}{t_{\mathrm{rec,RDCC}}}.
\end{equation}

Consequences:
\begin{itemize}
\item reionization begins later,
\item proceeds more gradually,
\item and ends at a similar redshift as in $\Lambda$CDM.
\end{itemize}

% ---------------------------------------------------------
\subsection{Reionization Redshift}

We define the reionization redshift $z_{\mathrm{re}}$ by
\begin{equation}
Q_{\mathrm{HII}}(z_{\mathrm{re}}) = 0.99.
\end{equation}

In RDCC, we find
\begin{equation}
z_{\mathrm{re,RDCC}}
\simeq
z_{\mathrm{re},\Lambda\mathrm{CDM}}
-
\Delta z,
\end{equation}
with
\begin{equation}
\Delta z \sim 0.3\text{--}0.6,
\end{equation}
depending on the relaxation scale $k_{\mathrm{rel}}$.

This shift is consistent with:
\begin{itemize}
\item the low optical depth measured by Planck,
\item the delayed star formation inferred from JWST,
\item the absence of strong early 21-cm absorption features.
\end{itemize}

% ---------------------------------------------------------
\subsection{Clumping Factor in RDCC}

The recombination rate depends on the clumping factor
\begin{equation}
C
=
\frac{\langle n_{\mathrm{H}}^{2} \rangle}{\langle n_{\mathrm{H}} \rangle^{2}}.
\end{equation}

In RDCC, suppressed small-scale structure reduces $C$:
\begin{equation}
C_{\mathrm{RDCC}}
=
C_{\Lambda\mathrm{CDM}}
\left[
1
-
\chi_{C}
\left(
\frac{k}{k_{\mathrm{rel}}}
\right)^{\alpha_{C}}
\right],
\end{equation}
with $\chi_{C} \sim 0.2$.

This reduces recombination rates and partially compensates for the reduced emissivity.

% ---------------------------------------------------------
\subsection{Thomson Optical Depth}

The Thomson optical depth is
\begin{equation}
\tau
=
c \sigma_{\mathrm{T}}
\int_{0}^{z_{\mathrm{max}}}
dz\,
\frac{n_{e}(z)}{(1+z)H(z)}.
\end{equation}

In RDCC:
\begin{equation}
\tau_{\mathrm{RDCC}}
=
\tau_{\Lambda\mathrm{CDM}}
\left[
1
-
\chi_{\tau}
\left(
\frac{k}{k_{\mathrm{rel}}}
\right)^{\alpha_{\tau}}
\right],
\end{equation}
with $\chi_{\tau} \sim 0.05$.

Typical values:
\begin{equation}
\tau_{\mathrm{RDCC}}
\simeq
0.050\text{--}0.055,
\end{equation}
consistent with Planck 2018.

% ---------------------------------------------------------
\subsection{Ionization History Summary}

RDCC predicts:
\begin{itemize}
\item a delayed onset of reionization,
\item a smoother ionization history,
\item reduced recombination rates due to lower clumping,
\item a low optical depth consistent with Planck,
\item and a characteristic scale dependence tied to $k_{\mathrm{rel}}$.
\end{itemize}

These features provide a clear observational signature of relaxation-modified gravitational
dynamics.

% ---------------------------------------------------------
% SECTION 3 — THERMAL EVOLUTION OF THE IGM IN RDCC
% ---------------------------------------------------------
\section{Thermal Evolution of the IGM in RDCC}

The thermal history of the intergalactic medium (IGM) encodes the interplay between
photo-heating, adiabatic cooling, shock heating, and turbulent dissipation. In the standard
$\Lambda$CDM model, the IGM temperature is primarily determined by the timing of
reionization and the subsequent balance between heating and cooling processes. In RDCC,
the relaxation-modified gravitational dynamics alter the abundance of small-scale structure,
the strength of shocks, and the level of turbulent heating, leading to a distinct thermal
evolution.

Relaxation affects the IGM through:
\begin{itemize}
\item reduced small-scale structure and weaker shocks,
\item delayed photo-heating due to reduced emissivity,
\item enhanced turbulent heating from relaxation-induced vorticity,
\item a flatter temperature–density relation,
\item a characteristic scale dependence tied to $k_{\mathrm{rel}}$.
\end{itemize}

% ---------------------------------------------------------
\subsection{Temperature–Density Relation}

The IGM temperature–density relation is commonly written as
\begin{equation}
T(\Delta,z)
=
T_{0}(z)\,
\Delta^{\gamma(z)-1},
\end{equation}
where $\Delta = \rho / \bar{\rho}$ is the overdensity.

In RDCC, both $T_{0}$ and $\gamma$ are modified:
\begin{itemize}
\item reduced small-scale structure lowers shock heating,
\item relaxation-induced vorticity enhances turbulent heating,
\item delayed reionization shifts the heating peak to lower redshift.
\end{itemize}

We write the RDCC temperature normalization as
\begin{equation}
T_{0,\mathrm{RDCC}}(z)
=
T_{0,\Lambda\mathrm{CDM}}(z)
\left[
1
-
\chi_{\mathrm{sh}}
+
\chi_{\mathrm{turb}}
\right],
\end{equation}
with $\chi_{\mathrm{sh}} \sim 0.05$ (reduced shocks) and $\chi_{\mathrm{turb}} \sim 0.02$ (enhanced turbulence).

The slope becomes
\begin{equation}
\gamma_{\mathrm{RDCC}}(z)
=
\gamma_{\Lambda\mathrm{CDM}}(z)
\left[
1
-
\chi_{\gamma}
\left(
\frac{k}{k_{\mathrm{rel}}}
\right)^{\alpha_{\gamma}}
\right],
\end{equation}
with $\chi_{\gamma} \sim 0.1$.

Consequences:
\begin{itemize}
\item slightly cooler low-density IGM,
\item flatter temperature–density relation,
\item broader temperature distribution.
\end{itemize}

% ---------------------------------------------------------
\subsection[Vorticity-Driven IGM Heating: The Lyman-alpha Turbulent Channel]
{Vorticity-Driven IGM Heating: The Lyman-$\alpha$ Turbulent Channel}

As shown in Companion~XVIII, relaxation-induced inter-sector vorticity acts as a non-thermal
heating channel for the low-density IGM. This mechanism injects turbulent energy into the
baryonic fluid on scales near the relaxation scale $k_{\mathrm{rel}}$, modifying the thermal history
and leaving a measurable imprint on the Lyman-$\alpha$ forest.

\subsubsection{Turbulent Heating Rate}

The vorticity-induced heating rate is
\begin{equation}
H_{\mathrm{vort}}(z)
=
\epsilon_{\mathrm{vort}}\,
\rho_{b}(z)\,
\exp\!\left[
-
\left(
\frac{k}{k_{\mathrm{rel}}}
\right)^{2}
\right],
\end{equation}
with $\epsilon_{\mathrm{vort}} \sim 0.01$–$0.03$.

This heating channel is strongest in underdense regions, where shock heating is weak.

\subsubsection{Modification of the Temperature–Density Relation}

The vorticity channel flattens the temperature–density relation:
\begin{equation}
\gamma_{\mathrm{RDCC}}(z)
=
\gamma_{\Lambda\mathrm{CDM}}(z)
\left[
1
-
\chi_{\mathrm{vort}}
\left(
\frac{k}{k_{\mathrm{rel}}}
\right)^{\alpha_{\mathrm{vort}}}
\right],
\end{equation}
with $\chi_{\mathrm{vort}} \sim 0.05$–$0.1$.

Consequences:
\begin{itemize}
\item reduced temperature contrast between voids and filaments,
\item broader temperature distribution,
\item improved agreement with Lyman-$\alpha$ forest constraints at $z \sim 4$–$6$.
\end{itemize}

\subsubsection[Lyman-alpha Forest Constraints]
{Lyman-$\alpha$ Forest Constraints}

The vorticity channel provides a CMB-independent lower bound on the relaxation scale:
\begin{equation}
k_{\mathrm{rel}}
\gtrsim
0.3\,h\,\mathrm{Mpc}^{-1}.
\end{equation}

This bound is complementary to:
\begin{itemize}
\item small-scale matter power spectrum constraints (Companion~XIX),
\item velocity-dispersion constraints (Companion~XXVII),
\item 21-cm heating constraints (Section~4).
\end{itemize}

% ---------------------------------------------------------
\subsection{Photo-Heating in RDCC}

Photo-heating during reionization is given by
\begin{equation}
H_{\mathrm{photo}}
=
\int_{\nu_{\mathrm{HI}}}^{\infty}
d\nu\,
J_{\nu}\,
\sigma_{\mathrm{HI}}(\nu)\,
(h\nu - h\nu_{\mathrm{HI}}).
\end{equation}

In RDCC:
\begin{itemize}
\item the ionizing background builds up more slowly,
\item the heating peak occurs later,
\item the total heating is slightly reduced.
\end{itemize}

We write
\begin{equation}
H_{\mathrm{photo,RDCC}}
=
H_{\mathrm{photo},\Lambda\mathrm{CDM}}
\left[
1
-
\chi_{\mathrm{ph}}
\left(
\frac{k}{k_{\mathrm{rel}}}
\right)^{\alpha_{\mathrm{ph}}}
\right],
\end{equation}
with $\chi_{\mathrm{ph}} \sim 0.05$.

% ---------------------------------------------------------
\subsection{Shock Heating in RDCC}

Shock heating scales with the velocity dispersion of collapsing regions. In RDCC:
\begin{itemize}
\item suppressed small-scale power reduces shock strengths,
\item delayed collapse shifts shock heating to lower redshift.
\end{itemize}

We model the RDCC shock heating rate as
\begin{equation}
H_{\mathrm{sh,RDCC}}
=
H_{\mathrm{sh},\Lambda\mathrm{CDM}}
\left[
1
-
\chi_{\mathrm{sh}}
\left(
\frac{k}{k_{\mathrm{rel}}}
\right)^{\alpha_{\mathrm{sh}}}
\right].
\end{equation}

% ---------------------------------------------------------
\subsection{Turbulent Heating in RDCC}

Relaxation-induced vorticity enhances turbulent energy injection into the IGM:
\begin{equation}
H_{\mathrm{turb,RDCC}}
=
H_{\mathrm{turb},\Lambda\mathrm{CDM}}
\left[
1
+
\chi_{\mathrm{turb}}
\left(
\frac{k}{k_{\mathrm{rel}}}
\right)^{\alpha_{\mathrm{turb}}}
\right],
\end{equation}
with $\chi_{\mathrm{turb}} \sim 0.02$.

Consequences:
\begin{itemize}
\item mild heating of low-density regions,
\item broader temperature distribution,
\item enhanced WHIM fraction.
\end{itemize}

% ---------------------------------------------------------
\subsection{Full RDCC Thermal Evolution}

The temperature evolution equation is
\begin{equation}
\frac{dT}{dt}
=
-2HT
+
\frac{2T}{3\Delta}
\frac{d\Delta}{dt}
+
\frac{2}{3k_{\mathrm{B}}n_{b}}
\left(
H_{\mathrm{photo}}
+
H_{\mathrm{sh}}
+
H_{\mathrm{turb}}
-
\Lambda_{\mathrm{cool}}
\right).
\end{equation}

In RDCC, each heating term is modified as above.

The resulting temperature evolution is:
\begin{itemize}
\item cooler IGM before reionization,
\item delayed heating peak,
\item slightly lower post-reionization temperature,
\item broader temperature distribution at all redshifts.
\end{itemize}

Typical values:
\begin{equation}
T_{0,\mathrm{RDCC}}(z=5)
\simeq
(7\text{--}9)\times 10^{3}\,\mathrm{K},
\end{equation}
consistent with Lyman-$\alpha$ forest constraints.

% ---------------------------------------------------------
\subsection{Summary}

RDCC modifies the thermal evolution of the IGM through:
\begin{itemize}
\item reduced shock heating,
\item enhanced turbulent heating,
\item delayed and reduced photo-heating,
\item flatter temperature–density relation,
\item and a characteristic scale dependence tied to $k_{\mathrm{rel}}$.
\end{itemize}

These effects produce a distinct thermal history that is testable with Lyman-$\alpha$ forest
data, 21-cm observations, and future IGM tomography.


% ---------------------------------------------------------
% SECTION 4 — MORPHOLOGY OF IONIZED REGIONS AND THE 21-CM SIGNAL IN RDCC
% ---------------------------------------------------------
\section{Morphology of Ionized Regions and the 21-cm Signal in RDCC}

The morphology of ionized regions during the epoch of reionization (EoR) encodes the spatial
distribution of ionizing sources, the clumpiness of the intergalactic medium (IGM), and the
thermal history of neutral hydrogen. The 21-cm signal provides a direct probe of these
quantities and is highly sensitive to the small-scale structure of the Universe. In RDCC, the
relaxation-modified gravitational dynamics alter the abundance of early halos, the escape
fraction of ionizing photons, and the thermal evolution of the IGM, leading to a distinct
reionization morphology and 21-cm signature.

Relaxation affects the 21-cm signal through:
\begin{itemize}
\item suppressed small-scale structure and reduced early emissivity,
\item enhanced escape fractions in low-mass halos,
\item reduced clumping and smoother ionization fronts,
\item delayed and weaker heating of the neutral IGM,
\item a characteristic turnover at the relaxation scale $k_{\mathrm{rel}}$.
\end{itemize}

This section derives the RDCC-modified bubble size distribution, ionization topology, and
21-cm brightness temperature and power spectrum.

% ---------------------------------------------------------
\subsection{Bubble Size Distribution}

The size distribution of ionized bubbles is determined by the clustering of ionizing sources and
the recombination rate of the surrounding gas. In RDCC:
\begin{itemize}
\item suppressed small-scale structure reduces the abundance of small bubbles,
\item enhanced escape fractions increase the growth of large bubbles,
\item reduced clumping factors lower recombination rates,
\item delayed reionization shifts bubble growth to lower redshift.
\end{itemize}

We model the RDCC bubble size distribution as
\begin{equation}
\frac{dn_{\mathrm{bub,RDCC}}}{dR}
=
\frac{dn_{\mathrm{bub},\Lambda\mathrm{CDM}}}{dR}
\left[
1
-
\chi_{\mathrm{small}}
\left(
\frac{R_{\mathrm{rel}}}{R}
\right)^{2}
+
\chi_{\mathrm{large}}
\left(
\frac{R}{R_{\mathrm{rel}}}
\right)^{\beta_{\mathrm{bub}}}
\right],
\end{equation}
with $\chi_{\mathrm{small}} \sim 0.2$ and $\chi_{\mathrm{large}} \sim 0.1$.

Consequences:
\begin{itemize}
\item fewer small ionized regions,
\item larger characteristic bubble sizes,
\item smoother ionization topology.
\end{itemize}

% ---------------------------------------------------------
\subsection{Ionization Topology}

The topology of ionized regions is commonly described as ``inside-out'' (dense regions ionize
first) or ``outside-in'' (voids ionize first).

In RDCC:
\begin{itemize}
\item reduced early star formation weakens the inside-out bias,
\item enhanced escape fractions allow photons to propagate further,
\item reduced clumping lowers recombination in voids,
\item turbulent heating smooths temperature gradients.
\end{itemize}

Thus, RDCC predicts a more homogeneous ionization topology with:
\begin{itemize}
\item larger ionized regions in voids,
\item weaker correlation between density and ionization,
\item smoother ionization fronts.
\end{itemize}

% ---------------------------------------------------------
\subsection{21-cm Brightness Temperature}

The differential 21-cm brightness temperature is
\begin{equation}
\delta T_{21}
\simeq
27\,\mathrm{mK}\,
x_{\mathrm{HI}}
(1+\delta)
\left(
1
-
\frac{T_{\gamma}}{T_{s}}
\right)
\left(
\frac{1+z}{10}
\right)^{1/2},
\end{equation}
where $T_{s}$ is the spin temperature.

In RDCC:
\begin{itemize}
\item delayed heating lowers $T_{s}$ at early times,
\item reduced shock heating lowers $T_{s}$ in overdensities,
\item enhanced turbulent heating raises $T_{s}$ in voids,
\item smoother ionization morphology reduces small-scale fluctuations.
\end{itemize}

The net effect is:
\begin{itemize}
\item a weaker global absorption trough,
\item a delayed transition to emission,
\item reduced small-scale structure in $\delta T_{21}$.
\end{itemize}

% ---------------------------------------------------------
\subsection{21-cm Power Spectrum}

The 21-cm power spectrum is defined as
\begin{equation}
\Delta^{2}_{21}(k)
=
\frac{k^{3}}{2\pi^{2}}\,P_{21}(k).
\end{equation}

In RDCC:
\begin{itemize}
\item suppressed small-scale structure reduces power at $k \gtrsim 0.5\,h\,\mathrm{Mpc}^{-1}$,
\item larger bubbles shift the peak to lower $k$,
\item smoother ionization fronts reduce the high-$k$ tail,
\item the relaxation scale $k_{\mathrm{rel}}$ imprints a characteristic turnover.
\end{itemize}

We write the RDCC power spectrum as
\begin{equation}
P_{21,\mathrm{RDCC}}(k)
=
P_{21,\Lambda\mathrm{CDM}}(k)
\left[
1
-
\chi_{21}
\left(
\frac{k}{k_{\mathrm{rel}}}
\right)^{2}
\exp\!\left(
-
\chi_{21,B}
\left(
\frac{k}{k_{\mathrm{rel}}}
\right)^{2}
\right)
\right],
\end{equation}
with $\chi_{21} \sim 0.1$ and $\chi_{21,B} \sim 0.05$.

Consequences:
\begin{itemize}
\item a smoother, more coherent 21-cm signal,
\item a characteristic suppression at small scales,
\item a turnover at $k \sim k_{\mathrm{rel}}$,
\item a clear observational signature for SKA and HERA.
\end{itemize}

% ---------------------------------------------------------
\subsection{Summary}

RDCC predicts:
\begin{itemize}
\item larger and more coherent ionized regions,
\item reduced small-scale structure in the ionization field,
\item a weaker and delayed 21-cm absorption trough,
\item a smoother 21-cm power spectrum with a turnover at $k_{\mathrm{rel}}$,
\item and a distinct ionization topology that differs from $\Lambda$CDM.
\end{itemize}

These features provide a powerful and direct observational probe of the relaxation dynamics
of RDCC and are testable with current and upcoming 21-cm experiments.

% ---------------------------------------------------------
% SECTION 5 — CONCLUSION
% ---------------------------------------------------------
\section{Conclusion}

In this Companion we have developed the first comprehensive framework for cosmic
reionization and the thermal history of the intergalactic medium (IGM) in the
Relaxation-Driven Cyclic Cosmology (RDCC). Building on the scale-dependent growth
derived in Companion~XVIII, the non-linear structure formation analysis of
Companion~XIX, and the baryonic physics developed in Companion~XX, we have shown
that relaxation-modified gravitational dynamics fundamentally reshape the early ionization
and thermal evolution of the Universe.

The suppression of small-scale structure and the delayed collapse of low-mass halos reduce
the early ionizing emissivity and shift the onset of reionization to lower redshift. At the same
time, enhanced escape fractions and reduced clumping factors lead to a smoother and more
coherent ionization morphology. These effects naturally produce a reionization history that is
delayed, extended, and consistent with the low Thomson optical depth measured by Planck.

The thermal evolution of the IGM is similarly modified. Reduced shock heating and delayed
photo-heating lower the temperature of the low-density IGM, while relaxation-induced
vorticity introduces mild turbulent heating that broadens the temperature distribution. The
resulting temperature–density relation is flatter and more homogeneous than in
$\Lambda$CDM, in agreement with Lyman-$\alpha$ forest constraints at $z \sim 5$–$6$.

The morphology of ionized regions reflects these changes. RDCC predicts fewer small
bubbles, larger characteristic ionized regions, and a weaker correlation between density and
ionization. This leads to a distinct 21-cm signature: a weaker and delayed global absorption
trough, reduced small-scale power, and a characteristic turnover in the 21-cm power spectrum
at the relaxation scale $k_{\mathrm{rel}}$. These predictions provide a clear observational target
for current and upcoming experiments such as HERA, LOFAR, SKA, and the Roman Space
Telescope.

Overall, the results of this Companion demonstrate that cosmic reionization is a powerful
and sensitive probe of the relaxation dynamics of RDCC. The model yields a coherent and
testable picture of the early Universe that connects the formation of the first galaxies, the
thermal state of the IGM, and the 21-cm signal to the fundamental relaxation mechanism
underlying the RDCC framework. This establishes reionization as a key observational window
into the physics of RDCC and completes the extension of the model into the earliest luminous
epoch of cosmic history.


\end{document}
